\subsection{Exact local data from the elementary spherical Hecke operators}

At a spherical place \(w\mid\ell\) apart from \(w_0,v'\), the
simultaneous Hecke congruence used in the odd-rank construction accepts
any collection of integral spherical operators away from \(\ell\).
Using the elementary symmetric operators in that collection gives a
second proof of the two exact local comparisons. The geometric seed
representations and the simultaneous congruence are the inputs to this
argument.

Write \(k=\overline{\mathbb Q}_\ell\) and \(\mathcal O_k=\{z\in k:|z|\le1\}\).
For a representation \(V\) of \(G_K\), with \(K/\mathbb Q_\ell\) finite,
write \(X(V)_{g,j}\) for the coefficient of \(t^j\) in
\(\det(1+tV(g))\), for \(g\in G_K\) and \(1\le j\le a=\dim V\).
These coefficients are integral on continuous representations, by a stable
lattice over a finite coefficient field as in Lemma~\ref{newgal:coeff}.

\begin{lemma}[Polynomial tests at a residual determinant]\label{addg:residual}
Let \((X_0,z_0)\) and \((X_i,z_i)\) consist of such coefficient systems
and an additional scalar in \(k\). Suppose for each \(\eta>0\) there
are finitely many \(i\) and coefficients \(c_i\in k\) such that
\[
 \left|P(X_0,z_0)-\sum_i c_iP(X_i,z_i)\right|\le\eta
 \qquad(P\text{ an integer polynomial in the coordinates}).
\]
Fix a finite coefficient DVR \(\mathcal O\subset k\) over which the
residual determinant \(\bar D\) of \(X_0\) is defined. Let \(R\) be
its universal local pseudodeformation ring over \(\mathcal O\) for
\(G_K\). For each \(\eta\), one can retain only the points with residual
determinant \(\bar D\), and choose weights \(d_i\) on these points,
such that their evaluation maps \(x_i\) and the target map \(x_0\) satisfy
\[
  \left|x_0(f)-\sum_i d_i x_i(f)\right|\le\eta\quad(f\in R),
  \qquad \left|z_0-\sum_i d_i z_i\right|\le\eta .
\]
\end{lemma}
\begin{proof}
The original tests hold for polynomials with coefficients in
\(\mathcal O_k\) at the same precision. Apply the integer test to each
monomial, multiply by its coefficient, and use the ultrametric inequality.
Choose a support and weights at precision \(\min(\eta,1/2)\).
For each point \(i\) with a
different residual determinant, some coordinate \(j_i\) differs from the
target modulo the maximal ideal of \(\mathcal O_k\). Indeed the
characteristic coefficients on group elements determine a determinant on
the group algebra by Amitsur's formula. Form the integral polynomial
\[
 H(X)=\prod_{i\,\mathrm{excluded}}
       \frac{X_{j_i}-(X_i)_{j_i}}{(X_0)_{j_i}-(X_i)_{j_i}} .
\]
The denominators are units. Testing \(HP\) gives the same inequality
for \(P\) with \(d_i=c_iH(X_i)\) on the retained points, since
\(H(X_0)=1\) and all excluded evaluations vanish. The constant test
ensures that some point remains.

The universal ring is topologically generated over \(\mathcal O\) by
the characteristic coefficients of group elements. The absolute Galois
group of a finite extension of \(\mathbb Q_\ell\) satisfies Mazur's
finiteness condition, so this ring is Noetherian and its topology is
the maximal-ideal topology. See Chenevier, Corollary 1.14,
Proposition 3.3, Remark 3.5, Example 3.6 and Proposition 3.7
of~\cite{chenevier}, with finite extension of the coefficient DVR.
The evaluations are continuous local maps into integer rings of finite
extensions. For an extension with larger residue field, the coefficient
values and \(\mathcal O\) lie in the inverse image of the chosen residue
field, so the same universal property applies. Approximate any \(f\in R\)
by integral coordinate polynomials in the maximal-ideal topology and
evaluate at the target and the retained points. There are finitely many
maps and the \(d_i\) are already fixed. Thus the precision at these maps
can be chosen to absorb \(\max(1,\max_i|d_i|)\). Passing to the limit
gives the first inequality. Testing the additional coordinate itself
gives the second.
\end{proof}

\begin{proposition}[Exact periods on the odd-rank spherical subfamily]
\label{addg:local}
Use the geometric seed and simultaneous congruence hypotheses of the
odd-rank construction in Sections~\ref{tr:section} and~\ref{newgal:section},
at a spherical \(w\mid\ell\) apart from both auxiliary places. For the
global representation \(r\) constructed from traces, \(V=r|_{G_{E_w}}\)
has the labelled jumps \(J_\tau\) of \eqref{newgal:jumps}. Its
Frobenius-semisimplified crystalline Weil parameter has the prescribed
Satake eigenvalues in the normalization of Proposition~\ref{newgal:phi}.
These conclusions hold for every coefficient prime \(\ell\).
\end{proposition}
\begin{proof}
Set \(K=E_w\). At a usable spherical prime \(v\nmid\ell\), put
\(q=\mathrm N v\) and let \(T_j\) be the hyperspecial double-coset
operator for \(j\) uniformizer entries. On spherical fixed vectors,
\(q^{\binom j2}T_j\) has eigenvalue the \(j\)-th elementary symmetric
function of \(q^{(a-1)/2}B_v\). The eigenvalue of \(T_j\) is
\(q^{j(a-j)/2}e_j(B_v)\), and the exponents
add to \(j(a-1)/2\). The representatives and coefficients are integral
at \(\ell\). A representation with zero fixed space satisfies the required
scalar-action assertion there as well.
\src{src/Development/AbelianVarietyTheory/Proofs/SatakeParameters/SphericalOperatorsExist.lean}{146}{the elementary spherical operators and their action}
and \src{src/Development/AbelianVarietyTheory/Statements/EllipticCurveLevels/HasLevel.lean}{36}{the admissible operator data}.

For a fixed Weil element \(s\) and power, let \(b_0\) denote the
predicted Weil trace and \(b_i\) its value at a seed. The congruence for
the singleton \(\{w\}\) includes the coordinate \(z=cb\), with one
nonzero scalar \(c\) fixed by the place, power and normalization. Take
as the away-from-\(\ell\) list all the preceding elementary operators on
usable spherical primes, indexed also by the prolongations and Frobenius
lifts as in Section~\ref{tr:section}. The congruence permits this list
and the local coordinate together, with the same coefficients for every
polynomial.
\src{src/Development/AbelianVarietyTheory/Definitions/GaloisGeometry/LimWU4OfLimWU3.lean}{46}{quantification over operator lists and the additional local coordinates}
and \src{src/Development/AutomorphicArithmetic/Foundations/GaloisRepresentations/TraceLimits/LimFinOddOf.lean}{185}{normalization of the Weil power}.

On a fixed finite support of this approximation the common good
geometric Frobenius elements are dense in the global Galois group.
The target identification away from \(\ell\) in Section~\ref{tr:section}
and the seed characteristic-polynomial identities identify all the
elementary coefficients at those elements. For a test in finitely many
group elements, approach each element by common good Frobenius elements,
keeping all its elementary degrees at the same prime. Hold \(z\) fixed.
Continuity of the representations passes the polynomial inequality to
the chosen elements, with the same finite support. Restriction to
\(G_K\) gives exactly the hypothesis of Lemma~\ref{addg:residual} for
\(X_0=X(V)\), \(X_i=X(V_i)\) and \(z_j=cb_j\). Taking the elementary
operators directly avoids division by factorials at the small primes.

Choose the fraction field of \(\mathcal O\) to contain \(c\) and the
images of all embeddings of \(K\) over \(\mathbb Q_\ell\).
Theorem~\ref{newgal:liu} first gives crystallinity of the actual
restriction \(V\), including its extension classes. All the local seeds
are crystalline of the common labelled type. In the ring \(R\) of the
lemma, take the semistable quotient \(R^{\mathbf v,1}\) with this Hodge
type and trivial inertial type, taking the semistability field to be
\(K\). The quotient and its characterization on field-valued points for
arbitrary residual multiplicities are given by A'Campo, Hevesi, Thorne
and Whitmore, Theorem 3.3.6 in
\href{https://arxiv.org/abs/2607.11763}{\emph{Local-global compatibility of automorphic Galois representations over CM fields at \(p\)}}.
The local pseudorepresentation records the semisimplification. Exactness
of crystalline periods and strictness of the de Rham filtration on
crystalline representations give the same labelled graded multiplicities
for a seed and its semisimplification, as also recorded for de Rham
representations in Lemma 3.3.19 of the same paper. In particular every retained
seed map kills the defining ideal. Apply Lemma~\ref{addg:residual} to
each element of this ideal and let \(\eta\) tend to zero. The target
factors through \(R^{\mathbf v,1}\). Applying the same exactness to the
already crystalline \(V\) gives its claimed labelled jumps.

The integers in each \(J_\tau\) are distinct by dominance, so this
type fits the highest-weight convention of \S3.3.9 of the same paper.
Use the spherical Bernstein block. By \S3.3.11 the invariant function
given by the trace at the prescribed Weil element and power is an element
\(F\in R^{\mathbf v,1}[1/\ell]\). Its specialization is the trace of
the semisimple Weil representation. This trace also agrees with the
trace on the original crystalline module by exactness. Choose \(h\ge0\)
and \(f\in R\) lifting \(\ell^h cF\). The choice is fixed
before \(\eta\). Each retained point has \(x_i(f)=\ell^h z_i\). The
two inequalities of the lemma then give
\[
       |x_0(f)-\ell^h z_0|\le\max(1,|\ell^h|)\eta .
\]
Letting \(\eta\) decrease to zero and cancelling \(\ell^h c\) gives
the predicted Weil trace on \(V\). Apply this to the first \(a\)
powers of geometric Frobenius and use Newton's identities in
characteristic zero. Crystallinity gives trivial inertia and zero
monodromy. This determines the asserted parameter.
\end{proof}

\subsection{Base change for the definite-unitary eigensystems}

\begin{proposition}[Attachment on the standard definite unitary group]
\label{addg:unitary}
Use the CM field, odd rank \(a\) and definite unitary group of
Section~\ref{newgal:section}, with the classical local Langlands
correspondence and Fontaine period functors. For each of its
automorphic representations \(U\) there is a continuous semisimple
\(r_U:G_E\to\mathrm{GL}_a(k)\). At every split spherical place
\(w\nmid\ell\) it is unramified with geometric Frobenius multiset
\(q_w^{(a-1)/2}B_w\). At every \(w\mid\ell\) its restriction is de Rham
with labelled jumps \(J_\tau\) of \eqref{newgal:jumps}. If \(U\) is
spherical there, it is crystalline and its crystalline Weil parameter
has geometric Frobenius multiset \(q_w^{(a-1)/2}B_w\).
This holds for all dominant integral weights and for any finite
exceptional set \(T\) of nonsplit places in the level.
\end{proposition}
\begin{proof}
The group in the construction is defined by \({}^{t}\!\bar g g=1\)
for the identity Hermitian form over \(E/E^+\). The field hypotheses
give an extension unramified at every finite place with every place
over \(\ell\) split. These are the standard assumptions of Newton and
Thorne, \S1, and their formula (1.0.1) is the displayed group.
Their map at a split place is projection to the chosen matrix factor,
which is the map defining the local component in \(B_w\). At infinity
the pullback of the algebraic module along the chosen embedding gives
the weight \(\lambda_{\iota\tau}\) used here. On the compact adelic
quotient the continuous automorphic functions of fixed finite level
and finite type at infinity are square integrable. Their finite-dimensional
translates at infinity are smooth by the representation theory of compact
Lie groups. An irreducible submodule of these automorphic forms therefore
occurs in the discrete automorphic spectrum. The projection to some
irreducible Hilbert summand is nonzero and maps this submodule into
its smooth, finite-type vectors.
\src{src/Development/AbelianVarietyTheory/Definitions/GaloisGeometry/LatticeDataThicknessBound.lean}{16}{field, automorphic and split local conditions}.
Theorem 1.2 and Corollary 1.3 of Newton and Thorne give the discrete
isobaric base change \(\Pi\) and \(r_U\), including compatibility
with \(\operatorname{rec}_{E_w}(\Pi_w|\det|^{(1-a)/2})\) at split places
and de Rhamness over \(\ell\). Their forward base change allows
ramification at nonsplit places. At a spherical place the parameter has
trivial inertia and zero monodromy. For a de Rham representation these
two conditions on its Weil--Deligne parameter give crystallinity by
semistable descent. The displayed norm twist gives the claimed factor
of \(q_w^{(a-1)/2}\).
Newton and Thorne,
\href{https://arxiv.org/abs/1912.11261}{\emph{Symmetric power functoriality for holomorphic modular forms}}, \S1.

For the labelled weights, let one discrete block of \(\Pi\) have
cuspidal constituent \(\mu\) on \(\mathrm{GL}_m\) and length \(b\).
Its archimedean constituents have twists \(\mu|\det|^{(b+1-2j)/2}\),
\(1\le j\le b\). In the construction of \(r_U\) this block contributes
\[
   \bigoplus_{j=1}^{b}
   r_{\mu|\det|^s}\otimes\varepsilon_\ell^{\,1-j},
   \qquad s=(m+b-1-a)/2.
\]
Here \(\mu|\det|^s\) is regular algebraic and polarizable. This is the
discrete-spectrum formula used in the proof of Newton and Thorne's
corollary. See Anastassiades and Thorne,
\href{https://arxiv.org/abs/1912.11267}{\emph{Raising the level of automorphic representations of
\(\mathrm{GL}_{2n}\) of unitary type}}, proof of Corollary 3.4.
For each labelled infinitesimal-character entry \(x\) of \(\mu\),
the cuspidal weight formula in Proposition~\ref{newgal:classical} and
the jump \(-1\) of \(\varepsilon_\ell\) give the entry
\[
  (m-1)/2-(x+s)+j-1
      =(a-1)/2-\bigl(x+(b+1-2j)/2\bigr).
\]
Taking all entries and blocks gives \((a-1)/2\) minus the infinitesimal
character of \(\Pi\), with multiplicity. By Theorem 1.2(3) this character
is that of the algebraic module of \(U\). Its entries in zero-based
indexing are \(\lambda_{\iota\tau,i}+(a-1)/2-i\). Subtraction gives
\(i-\lambda_{\iota\tau,i}\) for \(0\le i<a\), including equal adjacent
highest-weight entries.
\end{proof}

For example, the constant automorphic representation in odd rank
\(a\ge3\) has highest weight zero. Its algebraically normalized Satake multiset at a
split spherical place is \(\{1,q_w,\ldots,q_w^{a-1}\}\), so the attached
representation is \(\bigoplus_{i=0}^{a-1}\varepsilon_\ell^{-i}\) by
Chebotarev and semisimplicity. Its jumps \(0,\ldots,a-1\) fail
lacunarity since \(\{1\}\) and \(\{0,2\}\) have equal means.
For a general \(U\) the same uniqueness identifies \(r_U\) with any
semisimple trace limit matching its good Frobenius eigenvalues.
Proposition~\ref{addg:unitary} thus supplies the rank-\(a\) seed
and target Galois conclusions in the classical interpretation.
Realizing a particular seed in the prescribed cohomology tower,
proving the simultaneous automorphic congruence, and comparing any
separately specified period functors with Fontaine's are further assertions.

\subsection{Inserting the auxiliary cuspidal type at fixed weight}

\begin{proposition}\label{addg:auxiliary}
Fix any algebraic weight of the definite group and a level away from
the split auxiliary place \(w_0\mid\ell\). Let \(f\) be a simultaneous
eigenform of that weight and level, with system \(\theta\).
Take any list of integral spherical Hecke operators away from
\(\ell\) at the hyperspecial places of the level, or more generally
integral operators away from \(\ell\) acting by a scalar on each
relevant local fixed space. One may also include
fixed spherical double-coset operators or finite linear combinations
at hyperspecial places over \(\ell\) apart from \(w_0\), after a nonzero scalar
normalization for each operator. For every \(m\ge1\) there are finitely
many eigensystems \(\theta_i\) of the same weight and level away from
\(w_0\), all containing the auxiliary type of Section~\ref{newgal:section}
at \(w_0\), and coefficients \(c_i\in k\) such that
\[
   \left|P(\theta)-\sum_i c_iP(\theta_i)\right|
              \le |\ell|^m
      \qquad(P\text{ an integer polynomial in the listed operators}).
\]
\end{proposition}
\begin{proof}
In the finite-adelic model let \(A\) be the algebraic coefficient
representation, with \(f(\gamma g)=A(\gamma)f(g)\). Transport coefficients
by \(\iota^{-1}\), and measure the values of \(f\) using
\(A(g_\ell)^{-1}f(g)\). These values are left invariant. A lattice
preserved by the compact coefficient-place part of the original level
exists over a finite coefficient field, after shrinking the level at
\(w_0\) if needed. The resulting supremum norm is a maximum on the
finite double quotient and uses the same lattice at finer levels.
Integral Hecke operators away from \(\ell\) are contractive since their
finite sums of translations have trivial \(\ell\)-component. At each
marked coefficient place there are finitely many matrices \(A(t_\ell)\)
in a given operator. Multiplying its coefficients by a sufficiently
large power of \(\ell\) makes all the summands integral in the lattice.
These normalizations depend only on the operators and weight.
Rescale \(f\) so a coordinate evaluation \(L\) of norm at most one
satisfies \(L(f)=\|f\|=1\).

Inside a torsion-free compact open subgroup fixing \(f\) at \(w_0\),
choose an open \(J\) and a quotient
\(\theta_J:J\twoheadrightarrow C=\mathbb Z/\ell^m\mathbb Z\) such
that every nontrivial character of \(C\) inflated to \(J\) is a
supercuspidal type. Such subgroups exist for every prime by
Beuzart-Plessis, Theorem D.1 in the appendix to Fintzen and Shin,
\href{https://arxiv.org/abs/2009.08476}{\emph{Congruences of algebraic automorphic forms and supercuspidal representations}}.
For \(m=1\) one may compose a quotient of order \(\ell^2\)
with reduction modulo \(\ell\).
A sufficiently deep principal congruence subgroup is torsion-free by
the convergent matrix logarithm. For \(a>1\), surjectivity follows also
because a nontrivial character trivial on the image would occur in
the trivial representation. In rank one a coordinate of the logarithm
on the deep unit group gives a surjection to \(C\), and every smooth
irreducible representation is cuspidal. In level \(\ker\theta_J\), the right
action of \(C\) on the finite double quotient is free. Indeed the
rational stabilizer at level \(J\) is finite by compactness at infinity
and discreteness of the rational group in its adeles. Its projection
at \(w_0\), after conjugation, lies in the torsion-free group \(J\).
The faithful embedding of rational matrices into this completion
then makes the stabilizer trivial.

Choose one point of each \(C\)-orbit and let \(h\) be 1 minus
\(\ell^m\) times the characteristic function of the selected points.
Its sum on each orbit is zero and \(\|1-h\|=|\ell|^m\).
Since \(f\) is \(J\)-fixed, \(hf\) has only nontrivial
\(C\)-characters. Decomposing at the fixed weight and the finer
level gives a finite sum \(hf=\sum_i f_i\) of Hecke eigenvectors
in these character spaces. The decomposition follows from the
discrete automorphic decomposition and the scalar action of the
spherical operators on hyperspecial invariants. Each occurring
local representation at \(w_0\) contains the indicated type.
Every integral polynomial in the operators is contractive, so applying
\(L P(T)\) to \(f-hf\) proves the bound with \(c_i=L(f_i)\).
\end{proof}

This fixed-weight argument preserves any prescribed finite collection
of polynomial nonvanishing conditions already satisfied by the target.
Integral coefficients in \(k\) are allowed by the monomial argument of
Lemma~\ref{addg:residual}. After clearing denominators, let \(D\) be the product
of their integral test polynomials. Apply the estimate to \(DP\)
and divide by \(D(\theta)\ne0\). Terms with \(D(\theta_i)=0\)
drop out, and increasing \(m\) in advance absorbs the fixed loss
\(|D(\theta)|^{-1}\). For a polynomial required to be nonzero at
every ordering of a Satake multiset, take the product over permutations and express this
symmetric polynomial through the elementary spherical operators.
This proves the auxiliary-type step at arbitrary fixed weights,
including repeated highest-weight entries. Changing the weight to
produce lacunarity at \(v'\), or imposing a nonvanishing condition
that fails on the target, requires an additional argument.
